% ============================================================
%  Beamer presentation: Nonparametric Classification
%  Course of Machine Learning – Master Degree in CS
%  University of Rome "Tor Vergata" – A.Y. 2025-2026
%  Author: Giorgio Gambosi
% ============================================================

% !TEX root = all_beamer.tex
% ==============================================================================
% HEADER LATEX PER PRESENTAZIONI BEAMER
% Corso di Machine Learning - Giorgio Gambosi
% Università di Roma "Tor Vergata"
% ==============================================================================

% ------------------------------------------------------------------------------
% CLASSE DOCUMENTO E METADATI
% ------------------------------------------------------------------------------
\documentclass[aspectratio=169, 10pt]{beamer}

\subtitle{Course of Machine Learning}
\author{Giorgio Gambosi}
\institute[Tor Vergata]{Master's Degree in Computer Science\\University of Rome ``Tor Vergata''}
\date{A.Y. 2025--2026}

% ------------------------------------------------------------------------------
% PACCHETTI MATEMATICI
% ------------------------------------------------------------------------------
\usepackage{amsmath, amssymb, amsfonts, mathtools, amsthm}
\usepackage{bm}              % Simboli matematici in grassetto
\usepackage{upgreek}         % Lettere greche dritte
\usepackage{derivative}      % Gestione semplificata delle derivate

% ------------------------------------------------------------------------------
% PACCHETTI GRAFICI E VISUALIZZAZIONE
% ------------------------------------------------------------------------------
\usepackage{graphicx}        % Inclusione immagini
\usepackage{xcolor}          % Gestione colori
\usepackage{xkcdcolors}      % Colori aggiuntivi stile xkcd
\usepackage{microtype}       % Miglioramenti tipografici

% ------------------------------------------------------------------------------
% PACCHETTI PER TABELLE
% ------------------------------------------------------------------------------
\usepackage{booktabs}        % Tabelle professionali
\usepackage{array}           % Estensioni per tabelle
\usepackage{multirow}        % Celle multi-riga
\usepackage{multicol}        % Layout multi-colonna

% ------------------------------------------------------------------------------
% TIKZ E PGFPLOTS
% ------------------------------------------------------------------------------
\usepackage{tikz}
\usepackage{tikz-network}
\usepackage[customcolors]{hf-tikz}
\usepackage{pgfplots}
\pgfplotsset{compat=1.18}
\usepgfplotslibrary{fillbetween}

% Librerie TikZ
\usetikzlibrary{arrows, arrows.meta, decorations.pathmorphing, decorations.pathreplacing}
\usetikzlibrary{decorations.markings, snakes, positioning, backgrounds, fit, shadows}
\usetikzlibrary{automata, patterns, calc, intersections, through, shapes}
\usetikzlibrary{shapes.geometric, shapes.arrows, scopes, chains, matrix}

% ------------------------------------------------------------------------------
% ALTRI PACCHETTI
% ------------------------------------------------------------------------------
\usepackage[skins]{tcolorbox}  % Box colorati
\usepackage{subcaption}         % Sottocaption per figure
\usepackage{listings}           % Inclusione codice sorgente

% ------------------------------------------------------------------------------
% TEMA E STILE BEAMER
% ------------------------------------------------------------------------------
\usetheme{Madrid}
\usecolortheme{whale}
\useinnertheme{rounded}
\usefonttheme{professionalfonts}
\setbeamertemplate{navigation symbols}{}  % Rimuove simboli di navigazione

% ------------------------------------------------------------------------------
% DEFINIZIONE COLORI PERSONALIZZATI
% ------------------------------------------------------------------------------
% Colori principali del tema
\definecolor{MainBlue}{RGB}{31,56,100}      % Blu principale (header, titoli)
\definecolor{AccentBlue}{RGB}{70,130,180}   % Blu accento (blocchi)
\definecolor{AlertRed}{RGB}{160,30,30}      % Rosso per alert
\definecolor{ExGreen}{RGB}{30,120,60}       % Verde per esempi
\definecolor{BlockBg}{RGB}{220,232,246}     % Sfondo blocchi
\definecolor{torgray}{RGB}{220,225,235}     % Grigio Tor Vergata

% Colori alternativi (Tor Vergata style)
\definecolor{TvDark}{RGB}{30,50,80}         % Blu scuro header
\definecolor{TvMid}{RGB}{70,100,140}        % Blu medio blocchi
\definecolor{TvLight}{RGB}{210,220,235}     % Blu chiaro sfondi
\definecolor{TvAlert}{RGB}{160,30,30}       % Rosso alert
\definecolor{TvGreen}{RGB}{20,100,50}       % Verde esempi

% Colori supplementari
\definecolor{torred}{RGB}{153,0,0}
\definecolor{torcyan}{RGB}{0,112,192}
\definecolor{cerise}{HTML}{DA3A62}
\definecolor{teal}{HTML}{029386}
\definecolor{slateblue}{HTML}{5B5EA6}
\definecolor{orange}{HTML}{F97306}

% Colori per nodi TikZ
\definecolor{nodeblue}{RGB}{101,140,196}    % xkcd Faded Blue

% ------------------------------------------------------------------------------
% CONFIGURAZIONE COLORI BEAMER
% ------------------------------------------------------------------------------
\setbeamercolor{frametitle}{bg=MainBlue, fg=white}
\setbeamercolor{title}{bg=MainBlue, fg=white}
\setbeamercolor{structure}{fg=MainBlue}

% Blocchi standard
\setbeamercolor{block title}{bg=AccentBlue, fg=white}
\setbeamercolor{block body}{bg=BlockBg}

% Blocchi alert
\setbeamercolor{block title alerted}{bg=AlertRed, fg=white}
\setbeamercolor{block body alerted}{bg=red!8}

% Blocchi esempio
\setbeamercolor{block title example}{bg=ExGreen, fg=white}
\setbeamercolor{block body example}{bg=green!7}

% ------------------------------------------------------------------------------
% CONFIGURAZIONE FONT BEAMER
% ------------------------------------------------------------------------------
\setbeamerfont{frametitle}{size=\normalsize,series=\bfseries}
\setbeamerfont{block title}{size=\small,series=\bfseries}

% ------------------------------------------------------------------------------
% STILI TIKZ PERSONALIZZATI
% ------------------------------------------------------------------------------
\tikzset{
  >=latex,  % Frecce stile LaTeX di default
  % Nodi per reti neurali
  node/.style={thick,circle,draw=myblue,minimum size=22,inner sep=0.5,outer sep=0.6},
  node in/.style={node,green!20!black,draw=mygreen!30!black,fill=mygreen!25},
  node hidden/.style={node,blue!20!black,draw=myblue!30!black,fill=myblue!20},
  node convol/.style={node,orange!20!black,draw=myorange!30!black,fill=myorange!20},
  node out/.style={node,red!20!black,draw=myred!30!black,fill=myred!20},
  % Connessioni
  connect/.style={thick,mydarkblue},
  connect arrow/.style={-{Latex[length=4,width=3.5]},thick,mydarkblue,shorten <=0.5,shorten >=1},
  % Stili numerati per facile mappatura
  node 1/.style={node in},
  node 2/.style={node hidden},
  node 3/.style={node out}
}

% ------------------------------------------------------------------------------
% AMBIENTI TEOREMI
% ------------------------------------------------------------------------------
\setbeamertemplate{theorems}[numbered]
\newtheorem{mytheorem}{Theorem}
\newtheorem{mydefinition}{Definition}

% ------------------------------------------------------------------------------
% CONFIGURAZIONE LISTINGS (CODICE)
% ------------------------------------------------------------------------------
\lstset{
  language=Python,
  basicstyle=\ttfamily\tiny,
  keywordstyle=\color{MainBlue}\bfseries,
  commentstyle=\color{gray},
  frame=single,
  breaklines=true,
}

% ==============================================================================
% MACRO MATEMATICHE PERSONALIZZATE
% ==============================================================================

% ------------------------------------------------------------------------------
% SPAZI CALLIGRAFICI (Calligraphic spaces)
% ------------------------------------------------------------------------------
\providecommand{\calX}{\mathcal{X}}         % Spazio input
\providecommand{\calY}{\mathcal{Y}}         % Spazio output
\providecommand{\calH}{\mathcal{H}}         % Spazio ipotesi
\providecommand{\calT}{\mathcal{T}}         % Training set
\providecommand{\calA}{\mathcal{A}}         % Algoritmo
\providecommand{\calR}{\mathcal{R}}         % Rischio
\providecommand{\calF}{\mathcal{F}}         % Famiglia di funzioni
\providecommand{\calL}{\mathcal{L}}         % Loss function

% Alias per compatibilità
\providecommand{\Rr}{\mathcal{R}}
\providecommand{\Tr}{\mathcal{T}}
\providecommand{\Hr}{\mathcal{H}}

% ------------------------------------------------------------------------------
% RISCHIO E PROBABILITÀ
% ------------------------------------------------------------------------------
\providecommand{\barR}{\overline{\mathcal{R}}}  % Rischio empirico
\providecommand{\EmpR}{\overline{\mathcal{R}}}  % Rischio empirico (alias)
\providecommand{\Rbar}{\overline{\mathcal{R}}}  % Rischio empirico (alias)
\providecommand{\pM}{p_M}                       % Probabilità modello
\providecommand{\pC}{p_C}                       % Probabilità corretta
\providecommand{\Prob}[2]{\mathbb{P}_{#1}\!\left[#2\right]}  % Probabilità

% ------------------------------------------------------------------------------
% FUNZIONI INDICATRICI
% ------------------------------------------------------------------------------
\providecommand{\indicator}[1]{\mathbf{1}\!\left[#1\right]}  % Funzione indicatrice
\providecommand{\ind}[1]{\mathbf{1}[#1]}                      % Versione breve

% ------------------------------------------------------------------------------
% OPERATORI MATEMATICI
% ------------------------------------------------------------------------------
\providecommand{\argmax}{\operatorname*{arg\,max}}  % Argomento del massimo
\providecommand{\argmin}{\operatorname*{arg\,min}}  % Argomento del minimo
\providecommand{\vcd}[1]{\operatorname{VCdim}(#1)} % VC dimension
\providecommand{\defeq}{\stackrel{\text{def}}{=}}  % Definito come

% ------------------------------------------------------------------------------
% VETTORI E MATRICI (Bold Math)
% ------------------------------------------------------------------------------
% Vettori minuscoli
\providecommand{\bx}{\mathbf{x}}            % Vettore x
\providecommand{\bz}{\mathbf{z}}            % Vettore z
\providecommand{\bw}{\mathbf{w}}            % Vettore peso
\providecommand{\btt}{\mathbf{t}}           % Vettore target
\providecommand{\wvec}{\mathbf{w}}          % Vettore peso (alias)
\providecommand{\bth}{\boldsymbol{\theta}}  % Parametro theta

% Matrici maiuscole
\providecommand{\bX}{\mathbf{X}}            % Matrice dati
\providecommand{\bZ}{\mathbf{Z}}            % Matrice Z
\providecommand{\Xmat}{\mathbf{X}}          % Matrice dati (alias)
\providecommand{\Hmat}{\mathbf{H}}          % Matrice H
\providecommand{\Smat}{\mathbf{S}}          % Matrice scatter generica
\providecommand{\bS}{\mathbf{S}}            % Matrice scatter
\providecommand{\Ii}{\mathbf{I}}            % Matrice identità
\providecommand{\bZero}{\mathbf{0}}         % Vettore zero
\providecommand{\bPsi}{\boldsymbol{\Psi}}   % Matrice Psi

% Medie (overline)
\providecommand{\ox}{\overline{\mathbf{x}}} % Media vettore x
\providecommand{\ow}{\overline{\mathbf{w}}} % Media vettore w
\providecommand{\oX}{\overline{\mathbf{X}}} % Media matrice X

% Matrici scatter specifiche (analisi discriminante)
\providecommand{\SW}{\mathbf{S}_W}          % Within-class scatter
\providecommand{\SB}{\mathbf{S}_B}          % Between-class scatter
\providecommand{\ST}{\mathbf{S}_T}          % Total scatter
\providecommand{\GX}{G_{\mathbf{X}}}        % Matrice Gram

% ------------------------------------------------------------------------------
% COMANDI DI FORMATTAZIONE TESTO
% ------------------------------------------------------------------------------
\providecommand{\term}[1]{\textbf{\color{accent}#1}}        % Termine importante
\providecommand{\halert}[1]{{\color{AlertRed}\textbf{#1}}}  % Alert rosso
\providecommand{\mlalert}[1]{{\color{AccentBlue}\textbf{#1}}}  % Alert blu
\providecommand{\mlemph}[1]{{\color{MainBlue}\textit{#1}}}  % Enfasi blu

% ==============================================================================
% FINE HEADER
% ==============================================================================

\newcommand{\blankpage}{
\clearpage
\thispagestyle{empty}
\null
\clearpage
}


\definecolor{tvrblue}{RGB}{31,73,125}
\definecolor{tvrlightblue}{RGB}{173,198,224}
\definecolor{tvrdark}{RGB}{20,40,70}
\definecolor{tvrred}{RGB}{150,20,20}
\definecolor{tvrgreen}{RGB}{0,100,60}
\definecolor{tvrgray}{RGB}{245,245,245}

% Frame title
%\setbeamercolor{frametitle}{bg=tvrblue,fg=white}
%\setbeamercolor{framesubtitle}{fg=tvrlightblue}
%\setbeamerfont{frametitle}{size=\large,series=\bfseries}
%\setbeamerfont{framesubtitle}{size=\small,series=\normalfont}
%
%% Title page
%\setbeamercolor{title}{fg=white}
%\setbeamercolor{subtitle}{fg=tvrlightblue}
%\setbeamercolor{author}{fg=white}
%\setbeamercolor{institute}{fg=tvrlightblue}
%\setbeamercolor{date}{fg=tvrlightblue}
%\setbeamertemplate{title page}{%
%  \begin{tikzpicture}[remember picture,overlay]
%    \fill[tvrblue] (current page.north west) rectangle (current page.south east);
%    \node[anchor=center,text width=0.9\paperwidth,align=center]
%      at ([yshift=1.5cm]current page.center) {%
%        {\usebeamerfont{title}\usebeamercolor[fg]{title}\inserttitle}\\[0.5em]
%        {\usebeamerfont{subtitle}\usebeamercolor[fg]{subtitle}\insertsubtitle}};
%    \node[anchor=center,text width=0.9\paperwidth,align=center]
%      at ([yshift=-1.5cm]current page.center){%
%        {\usebeamerfont{author}\usebeamercolor[fg]{author}\insertauthor}\\[0.3em]
%        {\usebeamerfont{institute}\usebeamercolor[fg]{institute}\insertinstitute}\\[0.3em]
%        {\usebeamerfont{date}\usebeamercolor[fg]{date}\insertdate}};
%  \end{tikzpicture}%
%}
%
%% Header / footline
%\setbeamertemplate{headline}{}
%\setbeamertemplate{footline}{%
%  \leavevmode%
%  \hbox{%
%    \begin{beamercolorbox}[wd=.5\paperwidth,ht=2.4ex,dp=1.2ex,leftskip=4pt]{author in head/foot}%
%      {\tiny\color{white}\insertshortauthor\ (\insertshortinstitute)}%
%    \end{beamercolorbox}%
%    \begin{beamercolorbox}[wd=.35\paperwidth,ht=2.4ex,dp=1.2ex,center]{title in head/foot}%
%      {\tiny\color{white}\insertshorttitle}%
%    \end{beamercolorbox}%
%    \begin{beamercolorbox}[wd=.15\paperwidth,ht=2.4ex,dp=1.2ex,right,rightskip=4pt]{date in head/foot}%
%      {\tiny\color{white}\insertshortdate\quad\insertframenumber/\inserttotalframenumber}%
%    \end{beamercolorbox}}%
%  \vskip0pt%
%}
%\setbeamercolor{author in head/foot}{bg=tvrblue}
%\setbeamercolor{title in head/foot}{bg=tvrblue}
%\setbeamercolor{date in head/foot}{bg=tvrblue}
%
%% Navigation symbols off
%\setbeamertemplate{navigation symbols}{}
%
%% Block colours
%\setbeamercolor{block title}{bg=tvrblue,fg=white}
%\setbeamercolor{block body}{bg=tvrgray,fg=black}
%\setbeamercolor{block title alerted}{bg=tvrred,fg=white}
%\setbeamercolor{block body alerted}{bg=red!5,fg=black}
%\setbeamercolor{block title example}{bg=tvrgreen,fg=white}
%\setbeamercolor{block body example}{bg=green!5,fg=black}
%
%% Itemize bullets
%\setbeamertemplate{itemize item}{\color{tvrblue}$\bullet$}
%\setbeamertemplate{itemize subitem}{\color{tvrblue}$\circ$}
%\setbeamertemplate{enumerate item}{\color{tvrblue}\textbf{\insertenumlabel.}}
%
%% Section page (like outline slide)
%\AtBeginSection[]{%
%  \begin{frame}{Outline}
%    \tableofcontents[currentsection,hideothersubsections]
%  \end{frame}
%}
%
% ---- Packages ----------------------------------------------
\usepackage{amsmath,amssymb}
\usepackage{tikz}
\usetikzlibrary{arrows.meta,shapes,positioning,calc}
\usepackage{pgfplots}
\pgfplotsset{compat=1.18}
\usepackage{subcaption}
\usepackage{booktabs}
\usepackage{xcolor}
\usepackage{mathtools}

% ---- Macros ------------------------------------------------
\newcommand{\argmax}[1]{\underset{#1}{\operatorname{arg\,max}}}
\newcommand{\kh}{\kappa_h}

% ============================================================
\title{Nonparametric Classification}
\subtitle{Course of Machine Learning}
\author{Giorgio Gambosi}
\institute[Tor Vergata]{Master's Degree in Computer Science\\University of Rome ``Tor Vergata''}
\date{A.Y.\ 2025--2026}

% ============================================================
\begin{document}

% --- Title --------------------------------------------------
\begin{frame}[plain]
  \titlepage
\end{frame}

% --- TOC ----------------------------------------------------
\begin{frame}{Outline}
  \tableofcontents[hideallsubsections]
\end{frame}

% ============================================================
\section{Parametric vs.\ Nonparametric Approaches}
% ============================================================

\begin{frame}{Overview}
  \framesubtitle{Two families of probabilistic classifiers}
  \begin{block}{Core idea}
    Most probabilistic classifiers estimate (explicitly or implicitly)\\
    posterior class probabilities $p(C_k\mid\mathbf{x})$, then predict:
    \[
      \hat{C} = \argmax{k}\; p(C_k\mid\mathbf{x}).
    \]
  \end{block}
  \bigskip
  \begin{columns}[T]
    \column{0.47\textwidth}
      \textbf{\color{tvrblue}Parametric}
      \begin{itemize}
        \item Fixed functional form; finite parameter vector $\boldsymbol\theta$.
        \item Learn $\boldsymbol\theta$ from data.
        \item Compact, computationally efficient.
        \item Risk: misspecification $\Rightarrow$ biased predictions.
      \end{itemize}
    \column{0.47\textwidth}
      \textbf{\color{tvrblue}Nonparametric}
      \begin{itemize}
        \item Minimal shape assumptions.
        \item Model complexity grows with $n$.
        \item Flexible; consistent under mild conditions.
        \item Risk: curse of dimensionality; high compute cost.
      \end{itemize}
  \end{columns}
\end{frame}

% -----------------------------------------------------------
\begin{frame}{Parametric Approaches}
  \framesubtitle{Discriminative vs.\ generative}
  \begin{columns}[T]
    \column{0.47\textwidth}
      \begin{block}{Discriminative}
        Directly model $p(C_k\mid\mathbf{x};\boldsymbol\theta)$.\\[4pt]
        \emph{Example:} Softmax regression.\\[4pt]
        Training: maximize log-likelihood or minimize cross-entropy.
      \end{block}
    \column{0.47\textwidth}
      \begin{block}{Generative}
        Model class-conditional densities $p(\mathbf{x}\mid C_k;\boldsymbol\theta_k)$\\
        plus priors $p(C_k)$.\\[4pt]
        Posterior via Bayes:
        \[p(C_k\mid\mathbf{x})\propto p(\mathbf{x}\mid C_k)\,p(C_k).\]
        Also supports sampling, missing data.
      \end{block}
  \end{columns}
  \bigskip
  \textbf{Parameter estimation strategies:}
  \begin{itemize}
    \item \textbf{ML} – maximize likelihood of observed data.
    \item \textbf{MAP} – add prior over $\boldsymbol\theta$.
    \item \textbf{Bayesian} – maintain full posterior; average predictions.
  \end{itemize}
\end{frame}

% -----------------------------------------------------------
\begin{frame}{Nonparametric Approaches}
  \framesubtitle{Flexibility and its costs}
  \begin{block}{Key idea}
    Do \emph{not} commit to a fixed parametric family.\\
    Estimate the distribution \emph{directly} from data.\\
    Number of effective parameters increases with $n$.
  \end{block}
  \bigskip
  \begin{columns}[T]
    \column{0.47\textwidth}
      \textbf{\color{tvrgreen}Strengths}
      \begin{itemize}
        \item Adapt to arbitrary distributional shapes.
        \item Capture multimodality, skewness.
        \item Consistent as $n\to\infty$ (with proper tuning).
      \end{itemize}
    \column{0.47\textwidth}
      \textbf{\color{tvrred}Limitations}
      \begin{itemize}
        \item \textbf{Curse of dimensionality:} sparse neighborhoods in high $d$.
        \item High computational cost at prediction time.
        \item Sensitive to smoothing hyperparameters ($h$, $k$).
      \end{itemize}
  \end{columns}
\end{frame}

% ============================================================
\section{Kernel Density Estimation (Parzen Windows)}
% ============================================================

\begin{frame}{Motivation: Density from Local Counts}
  \framesubtitle{The fundamental identity}
  \begin{block}{Local counting principle}
    Let $\mathcal{R}(\mathbf{x})$ be a region around $\mathbf{x}$ with volume $V$.\\
    Probability mass in the region:
    \[
      P_\mathbf{x} = \int_{\mathcal{R}(\mathbf{x})} p(\mathbf{z})\,d\mathbf{z}.
    \]
  \end{block}
  \begin{itemize}
    \item With $n$ i.i.d.\ samples, let $k$ = number of samples in $\mathcal{R}(\mathbf{x})$.
    \item $k \sim \text{Binomial}(n, P_\mathbf{x})$, so $E[k] = nP_\mathbf{x}$.
    \item Empirical approximation: $P_\mathbf{x} \approx k/n$.
    \item If $p$ is approximately constant within $\mathcal{R}$: $P_\mathbf{x}\approx p(\mathbf{x})\cdot V$.
  \end{itemize}
  \begin{alertblock}{Fundamental density estimate}
    \[
      p(\mathbf{x}) \approx \frac{k}{nV}
    \]
    Density is large where many points fall into a \emph{small} region.
  \end{alertblock}
\end{frame}

% -----------------------------------------------------------
\begin{frame}{Two Strategies from the Same Identity}
  \framesubtitle{Volume fixed vs.\ count fixed}
  \[
    p(\mathbf{x}) \approx \frac{k}{nV}
  \]
  \bigskip
  \begin{columns}[T]
    \column{0.47\textwidth}
      \begin{block}{Fix $V$ $\Rightarrow$ Kernel Density Estimation}
        \begin{itemize}
          \item Set a region of fixed volume.
          \item Count how many samples fall in.
          \item $k$ varies with the data.
          \item Also called \emph{Parzen window} method.
        \end{itemize}
      \end{block}
    \column{0.47\textwidth}
      \begin{block}{Fix $k$ $\Rightarrow$ $k$-Nearest Neighbors}
        \begin{itemize}
          \item Fix the number of neighbors.
          \item Let the region expand until $k$ points are inside.
          \item $V$ varies with local data density.
          \item Automatically adapts to density.
        \end{itemize}
      \end{block}
  \end{columns}
\end{frame}

% -----------------------------------------------------------
\begin{frame}{KDE: Fixed-Volume Estimation}
  \framesubtitle{Parzen window with a hypercube kernel}
  \begin{itemize}
    \item Center a hypercube of side $h$ (volume $V = h^d$) at $\mathbf{x}$.
    \item Indicator kernel:
    \[
      k(\mathbf{z}) =
      \begin{cases}
        1 & \text{if } |z_i| \le \tfrac{1}{2},\ \forall\, i,\\
        0 & \text{otherwise.}
      \end{cases}
    \]
    \item Number of points inside: $\displaystyle k = \sum_{i=1}^n k\!\left(\frac{\mathbf{x}-\mathbf{x}_i}{h}\right)$.
  \end{itemize}
  \begin{block}{Parzen density estimate}
    \[
      p_n(\mathbf{x}) = \frac{1}{n h^d}\sum_{i=1}^{n} k\!\left(\frac{\mathbf{x}-\mathbf{x}_i}{h}\right)
    \]
  \end{block}
  \begin{itemize}
    \item Non-negative; integrates to 1.
    \item Interpretation: sum of ``bumps'' of mass $1/n$ centered at each data point.
  \end{itemize}
\end{frame}

% -----------------------------------------------------------
\begin{frame}{The Bandwidth Parameter $h$}
  \framesubtitle{Bias–variance trade-off in KDE}
  \begin{columns}[T]
    \column{0.47\textwidth}
      \begin{alertblock}{$h$ too large}
        \begin{itemize}
          \item Region contains many points.
          \item Estimate is \textbf{smooth} (low variance).
          \item May wash out structure: \textbf{high bias}.
        \end{itemize}
      \end{alertblock}
    \column{0.47\textwidth}
      \begin{alertblock}{$h$ too small}
        \begin{itemize}
          \item Region contains few points.
          \item Estimate follows random fluctuations.
          \item \textbf{High variance}, noisy.
        \end{itemize}
      \end{alertblock}
  \end{columns}
  \bigskip
  \begin{block}{Optimal $h$}
    Selecting $h$ = choosing a \textbf{bias–variance compromise}.\\
    Same tension as model complexity in supervised learning.
  \end{block}
\end{frame}

% -----------------------------------------------------------
\begin{frame}{Histogram Illustration of Bandwidth Effect}
  \framesubtitle{Varying $h$ changes the smoothness of the density estimate}
  % Reproduce the tikz histogram from the notes (simplified for beamer)
  \begin{center}
  \begin{tikzpicture}[scale=0.52]
    % Axes for h=2 (top)
    \draw[->,>=stealth](-5,8)--(5,8);
    \draw[->,>=stealth](-5,8)--(-5,9);
    % data dots
    \foreach \x in {-3.5,-3,-2.5,-2.5,-2.5,-2,-1.5,-1.5,-1,-1,-1,0,0,.5,.5,.5,.5,1,1,1.5,1.5,2,2,2,2.5,3.5,3.5,4,4,4.5}{
      \fill (\x,8.1) circle(0.04);}
    \node[left] at (-5.2,8.5) {\footnotesize $h=2$};

    % Axes for h=1 (middle)
    \draw[->,>=stealth](-5,5.5)--(5,5.5);
    \draw[->,>=stealth](-5,5.5)--(-5,6.5);
    \node[left] at (-5.2,6) {\footnotesize $h=1$};
    % Schematic histogram bars (h=1)
    \draw[red,thick](-5,5.5)--(-3,5.5)--(-3,5.7)--(-2,5.7)--(-2,5.9)--(-1,5.9)--(-1,6.1)--(0,6.1)--(0,6.3)--(1,6.3)--(1,6.1)--(2,6.1)--(2,5.9)--(3,5.9)--(3,5.7)--(4,5.7)--(4,5.5)--(5,5.5);

    % Axes for h=0.5 (bottom)
    \draw[->,>=stealth](-5,3)--(5,3);
    \draw[->,>=stealth](-5,3)--(-5,4);
    \node[left] at (-5.2,3.5) {\footnotesize $h=\varepsilon$};
    % Schematic spiky histogram (h=small)
    \draw[red,thick](-5,3)--(-3.5,3)--(-3.5,3.3)--(-3,3.3)--(-3,3)--(-2.5,3)--(-2.5,3.6)--(-2,3.6)--(-2,3.2)--(-1.5,3.2)--(-1.5,3.5)--(-1,3.5)--(-1,3)--(-0.5,3)--(-0.5,3.4)--(0,3.4)--(0,3.7)--(0.5,3.7)--(0.5,3.5)--(1,3.5)--(1,3.8)--(1.5,3.8)--(1.5,3.3)--(2,3.3)--(2,3.6)--(2.5,3.6)--(2.5,3)--(3,3)--(3,3.3)--(3.5,3.3)--(3.5,3.5)--(4,3.5)--(4,3.1)--(4.5,3.1)--(4.5,3)--(5,3);
  \end{tikzpicture}
  \end{center}
  \vspace{-0.3em}
  {\small
  \begin{itemize}
    \item \textbf{Top} ($h$ large): over-smoothed, structure washed out (high bias).
    \item \textbf{Bottom} ($h$ small): noisy, follows random fluctuations (high variance).
    \item Choosing the right $h$ is the central challenge of KDE.
  \end{itemize}}
\end{frame}

% -----------------------------------------------------------
\begin{frame}{Smooth Kernels and the Parzen Window}
  \framesubtitle{Beyond the histogram kernel}
  \begin{itemize}
    \item Histogram kernel: discontinuous, equal weight inside the region.
    \item \textbf{Smooth kernels} give \emph{higher} weight to nearby samples.
  \end{itemize}
  \medskip
  \begin{block}{Valid kernel requirements}
    $\displaystyle\int \kappa_h(\mathbf{u})\,d\mathbf{u} = 1$,\quad
    $\displaystyle\int \mathbf{u}\,\kappa_h(\mathbf{u})\,d\mathbf{u} = \mathbf{0}$,\quad
    $\displaystyle\int \|\mathbf{u}\|^2\kappa_h(\mathbf{u})\,d\mathbf{u} > 0$.
  \end{block}
  \medskip
  \begin{block}{Parzen window estimate (smooth kernel)}
    \[
      p(\mathbf{x}) = \frac{1}{n}\sum_{i=1}^{n}\kappa_h(\mathbf{x}-\mathbf{x}_i)
    \]
    Each training point contributes a ``bump'' of mass $1/n$.
  \end{block}
\end{frame}

% -----------------------------------------------------------
\begin{frame}{Common Kernel Functions}
  \framesubtitle{Box, Gaussian, and Epanechnikov}
  \begin{columns}[T,onlytextwidth]
    \column{0.32\textwidth}
      \begin{block}{Box / Histogram}
        \[k(\mathbf{u}) = \begin{cases}1 & |u_i|\le\tfrac{1}{2}\;\forall i\\0 & \text{o.w.}\end{cases}\]
        Discontinuous; equal weights.
      \end{block}
    \column{0.32\textwidth}
      \begin{block}{Gaussian}
        \[\kappa_h(\mathbf{u})=\frac{e^{-\|\mathbf{u}\|^2/2h^2}}{(\sqrt{2\pi}h)^d}\]
        Infinitely smooth; infinite support.
      \end{block}
    \column{0.32\textwidth}
      \begin{block}{Epanechnikov}
        \[\kappa_h(u)=\frac{3}{4h}\!\left(1-\frac{u^2}{h^2}\right)\]
        for $|u|\le h$, else $0$.\\
        Compact support; efficient.
      \end{block}
  \end{columns}
  \bigskip
  \begin{itemize}
    \item Choice of kernel affects \textbf{smoothness} and \textbf{computation}.
    \item Bandwidth $h$ governs the effective resolution regardless of kernel type.
    \item Gaussian and Epanechnikov are the most commonly used in practice.
  \end{itemize}
\end{frame}

% -----------------------------------------------------------
\begin{frame}{Kernel Shape Comparison}
  \framesubtitle{Visual comparison of common kernels for Parzen window estimation}
  \begin{columns}[c]
    \column{0.32\textwidth}
      \centering
      \begin{tikzpicture}[scale=0.65]
        \fill[red!40](-1,0)--(-1,1)--(1,1)--(1,0)--cycle;
        \draw[thick](-1,0)--(-1,1)--(1,1)--(1,0);
        \draw[very thin,gray](-1,-0.1) grid (1,1.1);
        \draw[->](-1.5,0)--(1.5,0) node[right]{\tiny$u$};
        \draw[->](0,0)--(0,1.5) node[above]{\tiny$k(u)$};
        \node at (0,-0.5){\footnotesize\textbf{Box}};
      \end{tikzpicture}
    \column{0.32\textwidth}
      \centering
      \begin{tikzpicture}[scale=0.65,domain=-1:1]
        \fill[green!40] plot[samples=50](\x,{1/(0.3*sqrt(2*3.14159))*exp(-\x*\x/(2*0.09))*0.28})--cycle;
        \draw[thick,samples=50,domain=-1:1] plot(\x,{1/(0.3*sqrt(2*3.14159))*exp(-\x*\x/(2*0.09))*0.28});
        \draw[very thin,gray](-1,-0.1) grid (1,1.1);
        \draw[->](-1.5,0)--(1.5,0) node[right]{\tiny$u$};
        \draw[->](0,0)--(0,1.5) node[above]{\tiny$\kappa(u)$};
        \node at (0,-0.5){\footnotesize\textbf{Gaussian}};
      \end{tikzpicture}
    \column{0.32\textwidth}
      \centering
      \begin{tikzpicture}[scale=0.65,domain=-1:1]
        \fill[blue!40] plot[samples=50](\x,{0.75*(1-\x*\x)})--cycle;
        \draw[thick,samples=50,domain=-1:1] plot(\x,{0.75*(1-\x*\x)});
        \draw[very thin,gray](-1,-0.1) grid (1,1.1);
        \draw[->](-1.5,0)--(1.5,0) node[right]{\tiny$u$};
        \draw[->](0,0)--(0,1.5) node[above]{\tiny$\kappa(u)$};
        \node at (0,-0.5){\footnotesize\textbf{Epanechnikov}};
      \end{tikzpicture}
  \end{columns}
  \vspace{0.6em}
  {\small Common kernel functions for Parzen window density estimation.\\
  Left: flat (box) kernel; Center: Gaussian kernel; Right: Epanechnikov kernel.}
\end{frame}

% -----------------------------------------------------------
\begin{frame}{Bandwidth Selection}
  \framesubtitle{Heuristics and cross-validation}
  \begin{block}{Bandwidth = central hyperparameter of KDE}
    Too large $\Rightarrow$ biased. \quad Too small $\Rightarrow$ unstable.
  \end{block}
  \medskip
  \textbf{Rule of thumb} (Gaussian reference):
  \[
    h^* = 1.06\,\hat{\sigma}\,n^{-1/5}
  \]
  where $\hat{\sigma}$ is the sample standard deviation.
  \medskip
  \begin{block}{Cross-validation (leave-one-out log-likelihood)}
    \[
      h^* = \argmax{h}\;\frac{1}{n}\sum_{i=1}^{n}\log p_{-i}(\mathbf{x}_i)
    \]
    $p_{-i}(\mathbf{x}_i)$: density estimate at $\mathbf{x}_i$ \emph{excluding} that point.\\
    Directly measures predictive performance on unseen data.
  \end{block}
\end{frame}

% -----------------------------------------------------------
\begin{frame}{Classification via KDE}
  \framesubtitle{Nonparametric generative classifier}
  \begin{itemize}
    \item Apply KDE \textbf{class-wise}: estimate $p(\mathbf{x}\mid C_i)$ for each class.
    \item Use only training points with label $C_i$:
    \[
      p(\mathbf{x}\mid C_i) = \frac{1}{n_i}\sum_{j:\,t_j=i}\kappa_h(\mathbf{x}-\mathbf{x}_j).
    \]
    \item Empirical class prior: $p(C_i) = n_i/n$.
  \end{itemize}
  \begin{block}{Decision rule (Bayes posterior)}
    \[
      \hat{C} = \argmax{i}\; p(\mathbf{x}\mid C_i)\,p(C_i)
    \]
  \end{block}
  \begin{itemize}
    \item Nonparametric analogue of \emph{generative classification}.
    \item No assumption on the shape of $p(\mathbf{x}\mid C_i)$.
    \item Number of ``parameters'' grows with $n$.
  \end{itemize}
\end{frame}

% ============================================================
\section{$k$-Nearest Neighbors}
% ============================================================

\begin{frame}{The $k$-NN Principle}
  \framesubtitle{Fixed count, variable volume}
  \begin{block}{Same starting point as KDE: $p(\mathbf{x})\approx k/(nV)$}
    But now: \textbf{fix} $k$, \textbf{let} $V$ adapt.
  \end{block}
  \begin{itemize}
    \item Let $r_k(\mathbf{x})$ = distance from $\mathbf{x}$ to its $k$-th nearest neighbor.
    \item Use a $d$-dimensional ball of radius $r_k(\mathbf{x})$: $V = c_d\,r_k(\mathbf{x})^d$.
    \item Density estimate:
    \[
      p(\mathbf{x}) \approx \frac{k}{n\cdot c_d\cdot r_k(\mathbf{x})^d}.
    \]
  \end{itemize}
  \begin{columns}[T]
    \column{0.47\textwidth}
      \textbf{Dense region:}\\
      $r_k(\mathbf{x})$ small $\Rightarrow$ high density estimate.
    \column{0.47\textwidth}
      \textbf{Sparse region:}\\
      $r_k(\mathbf{x})$ large $\Rightarrow$ low density estimate.
  \end{columns}
  \medskip
  \begin{alertblock}{}
    $k$-NN \textbf{automatically adapts} neighborhood to local data density.
  \end{alertblock}
\end{frame}

% -----------------------------------------------------------
\begin{frame}{$k$-NN Classification: Majority Vote Rule}
  \framesubtitle{Probabilistic derivation}
  \begin{itemize}
    \item Given query $\mathbf{x}$: find its $k$ nearest training points.
    \item Let $k_i$ = number of those $k$ neighbors in class $C_i$.
  \end{itemize}
  \begin{block}{Posterior via local counting}
    \[
      p(\mathbf{x}\mid C_i)=\frac{k_i}{n_i V},\qquad
      p(\mathbf{x})=\frac{k}{nV},\qquad
      p(C_i)=\frac{n_i}{n}.
    \]
    Bayes' rule gives:
    \[
      p(C_i\mid\mathbf{x})
      = \frac{p(\mathbf{x}\mid C_i)\,p(C_i)}{p(\mathbf{x})}
      = \frac{\tfrac{k_i}{n_i V}\cdot\tfrac{n_i}{n}}{\tfrac{k}{nV}}
      = \frac{k_i}{k}.
    \]
  \end{block}
  \begin{alertblock}{Decision rule: majority vote}
    \[
      \hat{C} = \argmax{i}\; k_i
    \]
    Posterior = fraction of neighbors in each class.
  \end{alertblock}
\end{frame}

% -----------------------------------------------------------
\begin{frame}{The Role of $k$: Smoothness Control}
  \framesubtitle{Small $k$ vs.\ large $k$}
  \begin{columns}[T]
    \column{0.47\textwidth}
      \begin{alertblock}{Small $k$ (e.g.\ $k=1$)}
        \begin{itemize}
          \item Highly local boundaries.
          \item Each training point = prototype.
          \item Risk of \textbf{overfitting}.
          \item High variance.
        \end{itemize}
      \end{alertblock}
    \column{0.47\textwidth}
      \begin{block}{Large $k$}
        \begin{itemize}
          \item Smoother decision boundary.
          \item Averages over more samples.
          \item Risk of \textbf{underfitting}.
          \item Higher bias.
        \end{itemize}
      \end{block}
  \end{columns}
  \bigskip
  \begin{block}{Choosing $k$ in practice}
    \begin{itemize}
      \item Cross-validation is the principled approach.
      \item Heuristic: $k \approx \sqrt{n}$.
      \item Optimal $k$ depends on noise level, class overlap, and $d$.
    \end{itemize}
  \end{block}
\end{frame}

% -----------------------------------------------------------
\begin{frame}{$k$-NN Decision Boundaries: Illustration}
  \framesubtitle{Effect of varying $k$ on the decision boundary (2-D dataset)}
  \begin{columns}[c]
    \column{0.32\textwidth}
      \centering
      \begin{tikzpicture}[scale=0.9]
        % Schematic: k=1 jagged boundary
        \fill[blue!20] (0,0) rectangle (2,2);
        \fill[red!20] (0,0) -- (0.5,0) -- (0.5,0.5) -- (1,0.5) -- (1,1) -- (1.5,1) -- (1.5,1.5) -- (2,1.5) -- (2,2) -- (0,2) -- cycle;
        \draw[thick,blue!70!black] (0.5,0) -- (0.5,0.5) -- (1,0.5) -- (1,1) -- (1.5,1) -- (1.5,1.5) -- (2,1.5);
        \draw[->](0,0)--(2.2,0) node[right]{\tiny$x_1$};
        \draw[->](0,0)--(0,2.2) node[above]{\tiny$x_2$};
        \node[below] at (1,-0.2) {\small $k=1$};
      \end{tikzpicture}
    \column{0.32\textwidth}
      \centering
      \begin{tikzpicture}[scale=0.9]
        % Schematic: k=5 smoother boundary
        \fill[blue!20] (0,0) rectangle (2,2);
        \fill[red!20] (0,0) -- (2,0) -- (2,1) .. controls (1.5,1.2) and (1,1.3) .. (0,2) -- cycle;
        \draw[thick,blue!70!black] (2,1) .. controls (1.5,1.2) and (1,1.3) .. (0,2);
        \draw[->](0,0)--(2.2,0) node[right]{\tiny$x_1$};
        \draw[->](0,0)--(0,2.2) node[above]{\tiny$x_2$};
        \node[below] at (1,-0.2) {\small $k=5$};
      \end{tikzpicture}
    \column{0.32\textwidth}
      \centering
      \begin{tikzpicture}[scale=0.9]
        % Schematic: k=15 very smooth boundary
        \fill[blue!20] (0,0) rectangle (2,2);
        \fill[red!20] (0,0) -- (2,0) -- (2,1.2) .. controls (1,1.4) .. (0,1.8) -- cycle;
        \draw[thick,blue!70!black] (2,1.2) .. controls (1,1.4) .. (0,1.8);
        \draw[->](0,0)--(2.2,0) node[right]{\tiny$x_1$};
        \draw[->](0,0)--(0,2.2) node[above]{\tiny$x_2$};
        \node[below] at (1,-0.2) {\small $k=15$};
      \end{tikzpicture}
  \end{columns}
  \vspace{0.5em}
  {\small Schematic of decision boundaries for a 2-D two-class problem.\\
  Left: $k=1$ gives jagged, overfit boundaries.\ Center: $k=5$ balances.\ Right: $k=15$ is smooth but may underfit.}
\end{frame}

% -----------------------------------------------------------
\begin{frame}{Theoretical Properties of $k$-NN}
  \framesubtitle{Asymptotic guarantees and finite-sample caveats}
  \begin{block}{Asymptotic consistency}
    With $k=k(n)$ chosen so that $k\to\infty$ and $k/n\to 0$ as $n\to\infty$:\\
    \[
      \text{Error}_{k\text{-NN}} \;\longrightarrow\; \text{Bayes error}\quad(n\to\infty).
    \]
    \emph{No assumption on the underlying distribution shape.}
  \end{block}
  \bigskip
  \textbf{Finite-sample considerations:}
  \begin{itemize}
    \item $k=1$: very local; overfits; each training point = prototype.
    \item Optimal $k$ depends on noise level, class overlap, and dimensionality.
    \item Cross-validation is the standard method for selecting $k$.
  \end{itemize}
\end{frame}

% -----------------------------------------------------------
\begin{frame}{Distance Metrics in $k$-NN}
  \framesubtitle{Critical design choice}
  \begin{alertblock}{$k$-NN is entirely distance-based}
    Feature scaling and metric choice can \textbf{dramatically} affect results.
  \end{alertblock}
  \medskip
  \textbf{Common distance metrics:}
  \begin{itemize}
    \item \textbf{Euclidean:} $\|\mathbf{x}-\mathbf{y}\|_2$ — default; affected by scale.
    \item \textbf{Manhattan:} $\|\mathbf{x}-\mathbf{y}\|_1$ — robust to outliers.
    \item \textbf{Mahalanobis:} accounts for feature correlations.
    \item \textbf{Cosine:} often used for text / high-dimensional sparse data.
  \end{itemize}
  \medskip
  \begin{block}{High-dimensional challenge}
    As $d$ increases, distances between points \textbf{concentrate}.\\
    All neighbors become equally distant $\Rightarrow$ notion of ``nearest'' weakens.\\
    One manifestation of the \textbf{curse of dimensionality}.
  \end{block}
\end{frame}

% -----------------------------------------------------------
\begin{frame}{Computational Aspects: Lazy Learning}
  \framesubtitle{Training vs.\ prediction cost}
  \begin{columns}[T]
    \column{0.47\textwidth}
      \begin{block}{Training phase}
        \begin{itemize}
          \item Minimal computation.
          \item Store the entire dataset.
          \item $O(nd)$ storage.
        \end{itemize}
        $k$-NN is a \textbf{lazy learner}: no explicit model is built.
      \end{block}
    \column{0.47\textwidth}
      \begin{block}{Prediction phase}
        \begin{itemize}
          \item Nearest-neighbor search required.
          \item Naive: $O(nd)$ per query.
          \item Acceleration: KD-trees, ball trees.
          \item Benefits diminish as $d$ grows.
        \end{itemize}
      \end{block}
  \end{columns}
  \bigskip
  \begin{itemize}
    \item All computation is \emph{deferred} to test time.
    \item Suitable when training data change frequently.
    \item Costly when $n$ and $d$ are both large.
  \end{itemize}
\end{frame}

% ============================================================
\section{KDE vs.\ $k$-NN: Comparison}
% ============================================================

\begin{frame}{KDE vs.\ $k$-NN: Side-by-Side}
  \framesubtitle{Two sides of the same local counting principle}
  \begin{center}
  \begin{tabular}{lcc}
    \toprule
    & \textbf{KDE / Parzen Window} & \textbf{$k$-NN} \\
    \midrule
    Fixed quantity & Volume $V$ & Count $k$ \\
    Varying quantity & Count $k$ & Volume $V$ \\
    Locality scale & Bandwidth $h$ & Neighborhood size \\
    Adaptation & Fixed scale everywhere & Adapts to local density \\
    Output & Smooth density & Posterior probability \\
    Continuity & Smooth (smooth kernel) & Piecewise constant \\
    \bottomrule
  \end{tabular}
  \end{center}
  \bigskip
  \begin{itemize}
    \item KDE: convenient when a natural scale $h$ is available or tunable.
    \item $k$-NN: advantageous in heterogeneous datasets with varying density.
    \item Both become more expressive as $n$ grows.
  \end{itemize}
\end{frame}

% -----------------------------------------------------------
\begin{frame}{Bias–Variance in KDE and $k$-NN}
  \framesubtitle{The same fundamental trade-off, different knobs}
  \begin{columns}[T]
    \column{0.47\textwidth}
      \begin{block}{KDE}
        \begin{itemize}
          \item $h$ large $\Rightarrow$ high bias, low variance.
          \item $h$ small $\Rightarrow$ low bias, high variance.
          \item Tuning: rule-of-thumb or CV.
        \end{itemize}
      \end{block}
    \column{0.47\textwidth}
      \begin{block}{$k$-NN}
        \begin{itemize}
          \item $k$ large $\Rightarrow$ high bias, low variance.
          \item $k$ small $\Rightarrow$ low bias, high variance.
          \item Tuning: cross-validation.
        \end{itemize}
      \end{block}
  \end{columns}
  \bigskip
  \begin{center}
  \begin{tikzpicture}[scale=0.85]
    \begin{axis}[
      xlabel={Complexity (small $h$ or small $k$)},
      ylabel={Error},
      xmin=0,xmax=10,ymin=0,ymax=1.2,
      width=8cm,height=4.5cm,
      xtick=\empty,ytick=\empty,
      axis lines=left,
      legend style={at={(0.97,0.97)},anchor=north east,font=\small}
    ]
      \addplot[blue,thick,domain=0.5:9.5,samples=50]{0.8*exp(-0.4*x)+0.05};
      \addlegendentry{Bias$^2$}
      \addplot[red,thick,domain=0.5:9.5,samples=50]{0.05+0.07*x};
      \addlegendentry{Variance}
      \addplot[black,thick,domain=0.5:9.5,samples=50]{0.8*exp(-0.4*x)+0.05+0.07*x};
      \addlegendentry{Total error}
    \end{axis}
  \end{tikzpicture}
  \end{center}
\end{frame}

% ============================================================
\section{Curse of Dimensionality}
% ============================================================

\begin{frame}{Curse of Dimensionality}
  \framesubtitle{Why locality fails in high dimensions}
  \begin{alertblock}{The fundamental problem}
    As $d$ increases, data become \textbf{sparse}:
    local neighborhoods either contain too few points or span an enormous fraction of the space.
  \end{alertblock}
  \bigskip
  \textbf{Concrete example:} To capture 10\% of the data volume in $d$ dimensions, one needs a hypercube of side:
  \[
    l = 0.1^{1/d}.
  \]
  \begin{center}
  \begin{tikzpicture}
    \begin{axis}[
      xlabel={Dimension $d$},
      ylabel={Side length $l$},
      xmin=1,xmax=20,ymin=0,ymax=1,
      width=7cm,height=4cm,
      axis lines=left,ytick={0,0.5,0.8,1},
      xtick={1,5,10,15,20},
      grid=major,grid style={dotted}
    ]
      \addplot[tvrblue,thick,domain=1:20,samples=50]{0.1^(1/x)};
    \end{axis}
  \end{tikzpicture}
  \end{center}
  {\small As $d$ grows, capturing a fixed fraction of data requires a neighborhood that spans almost the whole space.}
\end{frame}

% -----------------------------------------------------------
\begin{frame}{Manifestations of the Curse}
  \framesubtitle{Practical consequences for nonparametric methods}
  \begin{itemize}
    \item \textbf{Data sparsity:} exponentially more data needed to maintain density.
    \item \textbf{Distance concentration:} all pairwise distances become similar.
      \[\frac{\max_i d(\mathbf{x},\mathbf{x}_i) - \min_i d(\mathbf{x},\mathbf{x}_i)}{\min_i d(\mathbf{x},\mathbf{x}_i)} \to 0 \quad (d\to\infty).\]
    \item \textbf{Neighborhood size:} must grow to contain enough points, destroying locality.
    \item \textbf{Computational cost:} naive search cost grows with $n$ and $d$.
  \end{itemize}
  \medskip
  \begin{block}{Practical mitigations}
    \begin{itemize}
      \item Dimensionality reduction (PCA, autoencoders, \ldots).
      \item Feature selection — use only informative features.
      \item Metric learning — learn a task-adapted distance.
      \item Modern acceleration structures (FAISS, HNSW) for large-scale $k$-NN.
    \end{itemize}
  \end{block}
\end{frame}

% ============================================================
\section{Summary}
% ============================================================

\begin{frame}{Summary: Key Concepts}
  \framesubtitle{Nonparametric classification at a glance}
  \begin{columns}[T]
    \column{0.47\textwidth}
      \begin{block}{Kernel Density Estimation}
        \begin{itemize}
          \item Fix volume $V = h^d$.
          \item $p(\mathbf{x}) = \frac{1}{n}\sum_i \kappa_h(\mathbf{x}-\mathbf{x}_i)$.
          \item Bandwidth $h$: bias–variance knob.
          \item Extend to classification class-wise.
        \end{itemize}
      \end{block}
      \vspace{4pt}
      \begin{block}{$k$-NN Classification}
        \begin{itemize}
          \item Fix count $k$; let $V$ adapt.
          \item $p(C_i\mid\mathbf{x}) = k_i/k$.
          \item Decision: majority vote.
          \item $k$: bias–variance knob.
        \end{itemize}
      \end{block}
    \column{0.47\textwidth}
      \begin{block}{Shared Properties}
        \begin{itemize}
          \item Both implement a \emph{locality principle}.
          \item Both consistent as $n\to\infty$.
          \item Both suffer from the curse of dimensionality.
          \item Both require careful hyperparameter tuning.
        \end{itemize}
      \end{block}
      \vspace{4pt}
      \begin{alertblock}{Key difference}
        KDE: fixed scale, varying count.\\
        $k$-NN: fixed count, varying scale.
      \end{alertblock}
  \end{columns}
\end{frame}

% -----------------------------------------------------------
\begin{frame}{Key Formulas at a Glance}
  \framesubtitle{}
  \renewcommand{\arraystretch}{1.5}
  \begin{center}
  \begin{tabular}{ll}
    \toprule
    \textbf{Concept} & \textbf{Formula / Statement} \\
    \midrule
    Fundamental identity & $p(\mathbf{x}) \approx k/(nV)$ \\
    KDE estimate (smooth kernel) & $p(\mathbf{x}) = \frac{1}{n}\sum_i\kappa_h(\mathbf{x}-\mathbf{x}_i)$ \\
    Parzen (box kernel) & $p_n(\mathbf{x}) = \frac{1}{nh^d}\sum_i k\!\left(\frac{\mathbf{x}-\mathbf{x}_i}{h}\right)$ \\
    Rule-of-thumb bandwidth & $h^* = 1.06\,\hat\sigma\,n^{-1/5}$ \\
    $k$-NN density & $p(\mathbf{x})\approx k/(n c_d r_k(\mathbf{x})^d)$ \\
    $k$-NN posterior & $p(C_i\mid\mathbf{x}) = k_i/k$ \\
    KDE classification & $\hat C = \arg\max_i\; p(\mathbf{x}\mid C_i)\,p(C_i)$ \\
    \bottomrule
  \end{tabular}
  \end{center}
\end{frame}

%% -----------------------------------------------------------
%\begin{frame}{The Big Picture}
%  \framesubtitle{Where nonparametric methods fit}
%  \begin{center}
%  \begin{tikzpicture}[
%    node distance=1.4cm and 2.2cm,
%    box/.style={draw,rounded corners,fill=tvrgray,minimum width=2.8cm,minimum height=0.8cm,align=center,font=\small},
%    arrow/.style={->,>=stealth,thick}
%  ]
%    \node[box,fill=tvrblue!20](data){Training data\\$\{\mathbf{x}_i,t_i\}$};
%    \node[box,right=of data](local){Locality\\principle};
%    \node[box,above right=0.6cm and 1.8cm of local](kde){KDE\\(fix $V$)};
%    \node[box,below right=0.6cm and 1.8cm of local](knn){$k$-NN\\(fix $k$)};
%    \node[box,right=2.2cm of kde](kde_class){Class-wise\\density est.};
%    \node[box,right=2.2cm of knn](knn_class){Majority\\vote};
%    \node[box,fill=tvrgreen!30,right=1.8cm of kde_class,yshift=-0.7cm](dec){Decision\\$\hat{C}$};
%
%    \draw[arrow](data)--(local);
%    \draw[arrow](local)--(kde);
%    \draw[arrow](local)--(knn);
%    \draw[arrow](kde)--(kde_class);
%    \draw[arrow](knn)--(knn_class);
%    \draw[arrow](kde_class)--(dec);
%    \draw[arrow](knn_class)--(dec);
%  \end{tikzpicture}
%  \end{center}
%  \vspace{0.5em}
%  Both KDE and $k$-NN are \textbf{nonparametric}, \textbf{instance-based} classifiers that use
%  the locality principle to estimate class probabilities without assuming a parametric form.
%\end{frame}

\end{document}
